\documentclass[a4paper]{article}

\usepackage[fontset=fandol]{ctex}
\usepackage{amsmath, amssymb}
\usepackage{graphicx}
\usepackage[margin=2.5cm]{geometry}
\usepackage[most]{tcolorbox}
\usepackage{etoolbox}
\usepackage{listings}
\usepackage{hyperref}
\usepackage{booktabs}
\usepackage{subcaption}
\usepackage{float}
\usepackage{tikz}
\IfFileExists{pgfplots.sty}{
    \usepackage{pgfplots}
    \pgfplotsset{compat=1.18}
}{}

\newtcolorbox{knowledgebox}[1]{
    enhanced, colback=blue!5!white, colframe=blue!75!black, colbacktitle=blue!75!black,
    coltitle=white, fonttitle=\bfseries, title=#1, attach boxed title to top left={yshift=-2mm, xshift=2mm},
    boxrule=1pt, sharp corners
}

\newtcolorbox{importantbox}[1]{
    enhanced, colback=yellow!10!white, colframe=yellow!80!black, colbacktitle=yellow!80!black,
    coltitle=black, fonttitle=\bfseries, title=#1, sharp corners
}

\newtcolorbox{warningbox}[1]{
    enhanced, colback=red!5!white, colframe=red!75!black, colbacktitle=red!75!black,
    coltitle=white, fonttitle=\bfseries, title=#1, sharp corners
}

\lstset{
    language=Python,
    basicstyle=\ttfamily\small,
    keywordstyle=\color{blue},
    stringstyle=\color{red!60!black},
    commentstyle=\color{green!60!black},
    breaklines=true,
    frame=single,
    numbers=left,
    numberstyle=\tiny\color{gray},
    captionpos=b,
    extendedchars=false
}

\newcommand{\notetitle}{UCB CS294/194-196: Agentic AI (Fall 2025) - Lecture 03: Evolution of System Designs from an AI Engineer Perspective}
\newcommand{\noteauthors}{Yangqing Jia \& Codex}
\newcommand{\notedate}{2025-09-22}
\newcommand{\videochannel}{Agentic AI MOOC / UC Berkeley}
\newcommand{\videopublishdate}{2025-09-22}
\newcommand{\videoduration}{1:19:31}
\newcommand{\videourl}{https://www.youtube.com/watch?v=xqRAS6rAouo}
\newcommand{\videocoverpath}{../../raw/03_xqRAS6rAouo/03_xqRAS6rAouo.jpg}

\begin{document}

\begin{titlepage}
\centering
{\Large 课程笔记\par}
\vspace{1.2cm}
{\huge\bfseries \notetitle\par}
\vspace{0.8cm}
{\large \noteauthors\par}
\vspace{0.3cm}
{\large \notedate\par}
\vspace{1.2cm}

\ifdefempty{\videocoverpath}{
{\small 请在生成笔记时填入视频封面图的本地路径。\par}
}{
\includegraphics[width=0.82\textwidth,height=0.45\textheight,keepaspectratio]{\videocoverpath}\par
}

\vfill
\begin{tcolorbox}[width=0.9\textwidth, colback=black!2!white, colframe=black!60, sharp corners]
\textbf{视频作者/频道}：\videochannel\par
\textbf{发布日期}：\videopublishdate\par
\textbf{视频时长}：\videoduration\par
\textbf{视频链接}：
\ifdefempty{\videourl}{
未填写
}{
\href{\videourl}{\nolinkurl{\videourl}}
}
\end{tcolorbox}
\end{titlepage}

\tableofcontents
\newpage

\section{从研究者到工程师再到创业者}
Yangqing Jia 这一讲的叙事方式非常个人化：他先讲自己在 Berkeley、Google 和创业公司中的经历，再把这些经历抽象成一个工程判断。这个判断就是：\textbf{AI 系统的演化方式和传统软件系统不一样}。传统 CS 里，microservices、database、data analytics 的演化路径相对稳定；但 AI 模型、应用和基础设施同时在快速变化，因此 system design 也必须跟着变。

slide 3 用 ``Demystify LLM and AGI'' 开场，背后的含义并不是否定 AGI，而是提醒大家：很多看似神秘的系统，其实可以用非常朴素的工程原理解释。slide 4 的 ``Chinese Typewriter Problem'' 是一个极好的类比：中文输入从来不是简单的一键一字，而是通过一个间接编码/选择过程把想法映射到字符。这和今天的 agent / model interface 很像：表面上像“按键”，底层其实是 representation、encoding、routing 和 state management 的组合。

\begin{warningbox}{本讲的风格}
这是一场强烈带有产业经验和个人判断的讲座。它的价值在于观察一线系统如何演化，而不是把每句话当成普适定理。阅读时要抓住“趋势判断”而不是“字面语气”。
\end{warningbox}

\subsection{本章小结}
这一讲不是在讲单一技术点，而是在讲 AI 工程师如何观察产业：模型、应用和基础设施会共同变化，而系统设计必须顺着这个变化方向重构。

\section{新算法继续推动模型}
slide 8 到 11 的主线很明确：\textbf{新算法仍然在推动 LLM models}。讲者列举了 GPT-5、Grok、Gemini、Kimi、DeepSeek、Qwen、GPT-OSS 等例子，并强调“没有 bubble”的判断来自消费侧持续增长，而不只是训练侧的 hype。slide 11 把这条线讲得更抽象：GPT、MoE、test-time scaling、reinforcement learning 可以类比到计算机视觉和通用 AI 历史上的结构性创新。

这里最值得记住的，是他把 ``LLM'' 与 ``AGI'' 都做了去神秘化。对他来说，今天的 progress 并不来自玄学式的突然顿悟，而是来自多个简单原理的组合：更好的算法、更好的数据、更好的 scale、更好的 evaluation。正因为如此，他才会说历史会重复：AI 的每一轮提升，最终都会回到“方法论 + 计算”这个组合上。

\begin{knowledgebox}{这一讲对模型进展的判断}
\begin{itemize}
    \item 新算法仍然会继续推动模型前进；
    \item 关键变化不只是模型，而是模型与应用、基础设施的协同演化；
    \item ``看起来像 magic'' 的系统，往往可以用简单工程原则解释。
\end{itemize}
\end{knowledgebox}

值得注意的是，讲者并没有把重点放在“某个 benchmark 分数涨了多少”，而是放在产业演进曲线上。这种视角对 agent 课程很重要，因为 agent 不是纯研究问题，它天然牵涉模型迭代、系统接口和商业化反馈。

\begin{figure}[H]
\centering

\includegraphics[width=0.94\textwidth]{pdf_pages/page-11.png}
\caption{新算法持续推动 LLM 的典型轨迹：从 GPT、MoE 到 test-time scaling 与 RL。}
\end{figure}

\subsection{本章小结}
模型进步不是单点突破，而是算法、数据和算力三者持续叠加后的结果。理解这一点，才能理解为什么系统设计会被不断重写。

\section{应用层繁荣：consumer、prosumer 与 enterprise}
slides 13 到 19 转到应用层。这里的核心判断是：\textbf{模型和应用的耦合并没有想象中那么强}。slide 14 直接说，perfect app experience correlated, but independent from models。换句话说，模型能力很重要，但产品体验还取决于交互流程、任务分解、记忆、搜索、权限与 workflow。

slide 16 和 17 强调 consumer app landscape is highly fluid，说明消费级应用会随着 foundational models 的进步快速变化；而 slide 18 和 19 指出 prosumers 的 willingness to pay、enterprise 的慢变需求，才是更可能形成持续价值的区域。讲者把这总结为一个很实在的创业判断：如果你想做垂直领域的 agent，企业和专业用户往往比泛消费更容易形成稳定付费。

\begin{importantbox}{应用层的三个市场}
\begin{itemize}
    \item \textbf{Consumer}：增长快，但 landscape 非常流动；
    \item \textbf{Prosumer}：付费意愿更强，能驱动早期收入；
    \item \textbf{Enterprise}：需求更慢，但更适合把 domain knowledge 与新技术结合成可持续产品。
\end{itemize}
\end{importantbox}

这部分和 agent 课程的关系非常直接。所谓 agentic workflow，不是把一个聊天框套到每个任务上，而是把 memory、search、tools、human approval 和 business process 重新编排。对于真实产品来说，``模型能力'' 只是其中一个变量，``应用体验'' 才是最终收益函数。

\subsection{本章小结}
应用层的判断是：消费级应用会迅速变化，但 enterprise 与垂直场景更容易形成 agent 的真实价值。模型越强，应用层反而越需要被重新设计。

\section{Fall 2024 的系统前导层：frameworks、compound systems 与 workflow}
Berkeley Fall 2024 的前导课程对这一讲最有价值的补充，不是再讲一次“AI 会变快”，而是把 agent 系统里最常被口头化的几件事写成了显式工程对象：framework、workflow、compound system、trace 和 state transition。这样一来，\textbf{AI engineer perspective} 就不只是产业观察，而是可以落到系统结构上的设计语言。

`AutoGen` 把 multi-agent conversation 变成角色分工、消息路由和终止条件的组合；`StateFlow` 进一步告诉我们，很多 agent 任务需要的是 state-driven execution，而不是无约束对话。`Building a Multimodal Knowledge Assistant` 则把知识接入、图文索引和回答组织看成一条完整 pipeline，而不是单个模型能力。

`Compound AI Systems \& the DSPy Framework` 的贡献尤其明显：它把 prompt、demonstration、module composition 和 evaluator 变成可优化对象。对系统设计者来说，这比“写一个更长 prompt”重要得多，因为真正的 agent 方案往往是多个可替换模块的组合，而不是一个不可分解的魔法字符串。

`Agents for Software Development` 与 `AI Agents for Enterprise Workflows` 则把执行反馈写进真实作业系统：代码仓库、shell、测试日志、审批状态、企业表单和 trace 都成了 agent 要管理的环境状态。这个视角和本讲的产业判断是相通的：模型能力只是其中一层，真正决定能否落地的是 workflow 设计、状态管理和失败恢复。

\begin{importantbox}{Fall 2024 让这讲更像教材的地方}
\begin{itemize}
    \item frameworks 不是聊天壳，而是角色、状态与接口协议；
    \item compound systems 不是 prompt 堆叠，而是可优化的模块图；
    \item enterprise workflow 不是 demo，而是带 trace 和 failure recovery 的真实任务系统。
\end{itemize}
\end{importantbox}

\subsection{本章小结}
如果只看 Fall 2025，这一讲容易被读成“产业判断”；加上 Fall 2024 的前导层后，它更像一门教材里真正的系统设计章节：模型、应用和基础设施之间，必须通过 framework 与 workflow 连接起来。

\section{AI infra 是 IT 的第三根支柱}
slide 22 到 32 是这一讲最像系统课的部分。讲者把 AI infra 说成 IT strategy 的 third pillar，并且用历史类比解释它为什么和传统 cloud 不一样。传统路线是 scientific computing、web service cloud、data cloud，而 AI cloud 需要面对的是 exaflops 级算力、异构硬件、高带宽网络和训练/推理共存的资源需求。

\begin{figure}[H]
\centering
\includegraphics[width=0.94\textwidth]{pdf_pages/page-24.png}
\caption{AI infra 作为 IT 的第三根支柱：从 scientific computing、web service cloud、data cloud 到 AI cloud 的历史演化。}
\end{figure}

slide 25 的对比非常有启发性。Data compute 是 IO \(\gg\) compute，Web services 是 IO \(>\) compute，而 AI compute 是 compute \(\gg\) IO。也就是说，AI workload 在算力、分布式结构和工程复杂度上都不同于传统 web workload。它不是更大号的 web service，而是另一类系统。

slide 27 进一步解释 conventional cloud value proposition 为什么不再完全成立：传统云强调软件复杂性高、工作负载多样、供应链灵活和 VM 高可互换；而 AI cloud 里，软件更像 AI frameworks，workload 更统一，供应链灵活性更低，硬件和系统更像为单一类任务深度定制。这也是为什么讲者会说 ``AI infra is different, but also the same''：很多云工程原理仍然有效，但系统边界已经变了。

\begin{knowledgebox}{AI infra 的新要求}
\begin{itemize}
    \item 面向模型训练与推理的高带宽、高吞吐和低延迟；
    \item 面向开发者效率的统一平台；
    \item 面向模型生命周期的 dev / train / inference 一体化；
    \item 面向供应链和资源调度的多云、弹性与利用率管理。
\end{itemize}
\end{knowledgebox}

他对 baremetal 和 K8s 的那句评论不应当机械理解，而应当读成一个系统判断：通用云栈并不能自动满足 AI 工作负载的调度、监控、恢复和效率要求。要做 AI infra，需要更贴近训练、推理和模型生命周期的基础设施抽象。

\subsection{本章小结}
AI infra 不是把传统云换个名字，而是围绕 AI workload 重新设计平台。它之所以成为第三根支柱，是因为 AI 已经有了自己的资源结构、调度结构和商业结构。

\section{硬件与软件设计：AI 时代的 back to the future}
slide 30 到 32 把讲座收束到系统设计层：what you want to care about 是 developer efficiency 和 infra efficiency。开发者时间宝贵，GPU 也会故障，而且比很多人想的更频繁，所以真正可用的系统必须从开发、训练到推理都能稳定运行。

讲者在这里提出了一个很实用的 best practices 清单：multi-cloud supply chain management、elasticity and utilization management、AI-native platform，以及围绕 model and applications 组建自己的团队。它们都指向同一个事实：AI 系统最难的部分，往往不是算法本身，而是让算法在真实约束里持续工作。

\begin{figure}[H]
\centering
\includegraphics[width=0.94\textwidth]{pdf_pages/page-27.png}
\caption{传统云与 AI 云的差异：软件复杂度、工作负载一致性和供应链属性都发生了变化。}
\end{figure}

后半段 transcript 还明确提到了更有前瞻性的趋势：一方面，机器人、手机智能和 edge compute 会推动更分布式的 AI 资源部署；另一方面，中心化的大规模训练和推理集群仍然会存在。于是系统设计的真实问题变成了：如何在 centralized scaling 和 edge deployment 之间取得平衡。

\begin{importantbox}{讲者的核心工程判断}
今天的 AI 系统设计，是在两个方向之间找平衡：
\begin{itemize}
    \item 一边是更大、更集中、更高效的训练和推理集群；
    \item 另一边是更靠近设备、更贴近场景、更节能的 edge / robotics / mobile intelligence。
\end{itemize}
这不是二选一，而是同一产业的两头。
\end{importantbox}

他在最后把这个平衡总结为 exploration 与 exploitation：当大家还在争夺最佳模型时，愿意多烧一些资源；而当工程问题已经清晰时，efficient models、distributed compute 和更贴近部署位置的资源配置会越来越重要。robotics 之所以会再次变得重要，也正是因为它把这条趋势推到了物理世界。

\subsection{本章小结}
AI 时代的硬件与软件设计，不再只是“更快的服务器”，而是围绕模型、场景和部署位置重新做平衡。中心化算力和 edge 智能会同时增长，这就是讲者说的 ``back to the future''。

\section{总结与延伸}
这一讲的真正主题，是从 AI engineer 的视角重新看 system design 的演化。它告诉我们：
\begin{enumerate}
    \item 新算法仍在推动模型，但模型不再是唯一中心；
    \item 应用层繁荣，尤其是 enterprise 和 vertical tasks，决定了 agent 的商业空间；
    \item AI infra 已经成为独立的 IT 支柱；
    \item 面向 AI 的硬件与软件设计，正在回到“算力、带宽、调度与部署”这些老问题，只是问题规模更大了。
\end{enumerate}

对 agent 课程而言，这一讲提供了一个重要底色：\emph{agents 不是只靠模型能力就能落地的，它们需要完整的系统生态。} 也就是说，如果你只看 benchmark 分数，会错过真正决定产品成败的基础设施层。

\subsection{本章小结}
Lecture 03 的价值在于把产业、应用和基础设施连成一条线。它不是讲某个单点技术，而是在讲 AI 工程师如何判断：哪些问题会被模型改写，哪些问题会被系统重塑。

\end{document}
